% !TeX spellcheck = en_US

%%%%%%%%%%%%%%%%%%%%%%%
\begin{figure}
	\centering
	\begin{tabular}{c}
		\bottomAlignProof
		\AxiomC{$ $}
		\RightLabel{\scriptsize \textsc{skip}}
		\UnaryInfC{$\langle \lstiC{SKIP}, \sigma \rangle \Downarrow \sigma$}
		\DisplayProof
		\qquad
		\bottomAlignProof
		\AxiomC{$\sigma(\lstiC{x}) = (v, s)$}
		\RightLabel{\scriptsize \textsc{inc}}
		\UnaryInfC{$\langle \lstiC{INC x}, \sigma \rangle \Downarrow \sigma[\lstiC{x} \to (v+1,s)]$}
		\DisplayProof
		\\ \\
		\bottomAlignProof
		\AxiomC{$\sigma(\lstiC{x}) = (v, s)$}
		\RightLabel{\scriptsize \textsc{dec}}
		\UnaryInfC{$\langle \lstiC{DEC x}, \sigma \rangle \Downarrow \sigma[\lstiC{x} \to (v-1,s)]$}
		\DisplayProof
		\qquad
		\bottomAlignProof
		\AxiomC{$\sigma(\lstiC{x}) = (v, s)$}
		\RightLabel{\scriptsize \textsc{push}}
		\UnaryInfC{$\langle \lstiC{PUSH x}, \sigma \rangle \Downarrow \sigma[\lstiC{x} \to (0,v::s)]$}
		\DisplayProof
		\\ \\
		\bottomAlignProof
		\AxiomC{$\sigma(\lstiC{x}) = (v,
        s)$}
		\RightLabel{\scriptsize \textsc{pop}}
		\UnaryInfC{$\langle \lstiC{POP x}, \sigma \rangle \Downarrow \sigma[\lstiC{x} \to (\hd(s),\tl(s))]$}
		\DisplayProof
		\qquad
		\bottomAlignProof
		\AxiomC{$\langle \lstiC{P} , \sigma \rangle \Downarrow \nu $}
		\AxiomC{$\langle \lstiC{Q} , \nu \rangle \Downarrow \tau $}
		\RightLabel{\scriptsize \textsc{seq}}
		\BinaryInfC{$\langle \lstiC{P;Q}, \sigma \rangle \Downarrow \tau$}
		\DisplayProof
        \\ \\
%        ---------------- iteration
		\bottomAlignProof
		\AxiomC{$\sigma(\lstiC{x}) = (v,s)$}
		\AxiomC{$v \geq 0$}
		\AxiomC{$\langle \lstiC{P}, \sigma \rangle \Downarrow^v \tau$}
		\RightLabel{\scriptsize \textsc{for-fwd}}
		\TrinaryInfC{$\langle \lstiC{FOR x \{P\}}, \sigma \rangle \Downarrow \tau$}
		\DisplayProof
		\\ \\
		\bottomAlignProof
		\AxiomC{$\sigma(\lstiC{x}) = (v,s)$}
		\AxiomC{$v < 0$}
		\AxiomC{$
         \langle \lstiC{-(P)}, \sigma \rangle \Downarrow^{-v} \tau$}
		\RightLabel{\scriptsize \textsc{for-bwd}}
		\TrinaryInfC{$\langle \lstiC{FOR x \{P\}}, \sigma \rangle \Downarrow \tau$}
		\DisplayProof
		\\ \\
		\bottomAlignProof
		\AxiomC{$ $}
		\RightLabel{\scriptsize \textsc{base}}
		\UnaryInfC{$\langle \lstiC{P}, \sigma \rangle \Downarrow^0 \sigma$}
		\DisplayProof
		\qquad
		\bottomAlignProof
		\AxiomC{$\langle \lstiC{P}, \sigma \rangle \Downarrow \nu $}
		\AxiomC{$\langle \lstiC{P}, \nu \rangle \Downarrow^{n} \tau$}
		\RightLabel{\scriptsize \textsc{step}}
		\BinaryInfC{$\langle \lstiC{P}, \sigma \rangle \Downarrow^{n+1} \tau$}
		\DisplayProof
	\end{tabular}
	\caption{\nsemantics of \SCORE, the naive one.}
	\label{fig:scoreSemanticnaive}
\end{figure}
%%%%%%%%%%%%%%%%%%%%%%%

%-------------
\section{A classical non-reversible semantics for \SCORE to start with}
\label{section:A naive semantic ldots irreversible}

We here introduce \nsemantics  for \SCORE, its `\textsf{N}' standing for `naive', word that emphasizes that the choices to define \nsemantics are simplistic, even though they provide a starting point for the other semantics introduced later.

\vspace{\baselineskip}\noindent
The set $\StackZZ$ of all and only the stacks containing values of $\ZZ$ is inductively given as:
\begin{itemize}
\item
$[\,] \in \StackZZ$, and
\item
$h \cons t \in \StackZZ$ if $h \in \mathbb{Z}$ and $t \in \StackZZ$
\end{itemize}
and is associated with the following head and tail functions:
\begin{align}
\label{align: hd and tl base}
\hd(\nil)     &= 0  &\tl(\nil)     &= \nil \\
\label{align: hd and tl ind}
\hd(h\cons t) &= h  &\tl(h \cons t) &= t
\enspace ,
\end{align}
\noindent
which are both \emph{total}, coherently with our interest about \TIF.

\begin{definition}[\nsemantics]
\label{definition:States of nsemantics}
Fig.~\ref{fig:scoreSemanticnaive} shows the rules of \nsemantics which operates on states $\sigma : V \to \ZZ \times \StackZZ$ such that, if $\sigma(\lstiC{x}) = (v, s)$ then:
\begin{itemize}
\item
$v$ is the value currently hold by \lstiC{x} in $\sigma$;
\item
$s$ is the stack currently hold by \lstiC{x} in $\sigma$, containing all and only the values previously stored in \lstiC{x} immediately before every interpretation of ``\lstiC{PUSH x}'' which resets \lstiC{x} to $0$.
\end{itemize}
\noindent
States are ranged over by $\Sigma$.
As usual, $\sigma[\lstiC{x} \to (v', s')]$ is $\sigma$ where \lstiC{x} maps to $(v', s')$, for some $v', s'$, and every \lstiC{y} other than \lstiC{x} maps to $\sigma(\lstiC{y})$.
\end{definition}
\noindent
The rule \textsc{push} interprets ``\lstiC{PUSH x}'' by pushing $v$ on top of $s$ and replacing $v$ by $0$, for any $\sigma$ and $x$ such that $\sigma(x) = (v, s)$. 
The rule \textsc{pop} interprets ``\lstiC{POP x}'' by popping the top $h$ of $s$ and using $h$ in place of $v$.


%%%%------------------
\subsection{\nsemantics always terminates}
The reason why, for every \lstiC{P}, and state $\sigma$, the interpretation of \lstiC{P} in $\sigma$ by means of \nsemantics always terminates in some unique state $\tau$ is almost straightforward.
Every iteration of ``\lstiC{FOR x \{P\}}'' generates the interpretation of a finite sequence ``\lstiC{P;...;P}'', regardless of whether ``\lstiC{FOR x \{P\}}'' is interpreted using \textsc{for-fwd} or \textsc{for-bwd}, as both unfold into a tree of multiple \textsc{step} occurrences with a single occurrence of \textsc{base} as its leaf.
Both ``\lstiC{INC x}'' and ``\lstiC{DEC x}'' rely on $\ZZ$ being infinite in both directions, meaning that there is no inherent limit to the number of times they can be applied.
The same observation holds for ``\lstiC{PUSH x}'' and ``\lstiC{POP x}''. Even when $\sigma(\lstiC{x}) = (v, \nil)$, ``\lstiC{POP x}'' can be applied indefinitely because, based on \eqref{align: hd and tl base}, $\nil$ is a fixed point of its interpretation $\tl(\cdot)$.

%%%%------------------
\subsection{\nsemantics is not reversible}
\label{subsection:nsemantics is not reversible}

This is because, for every \lstiC{x}, at least one state $\sigma$ exists such that \LH{($\langle \lstiC{POP x;} \invSC{(POP x)}, \sigma \rangle \not\Downarrow \sigma$)}, according to the rules in Fig.~\ref{fig:scoreSemanticnaive}.

For example, let $\sigma$ be such that $\sigma(\lstiC{x}) = (5,[2])$. 
Then, $\langle \lstiC{POP x}, \sigma \rangle \Downarrow \sigma[\lstiC{x} \to (2, \nil)]$. As $\lstiC{PUSH x} =$  $\invSC{(POP x)}$, we have $\langle \invSC{(POP x)}, \sigma[\lstiC{x} \to (2, \nil)] \rangle \Downarrow \sigma[\lstiC{x} \to (0,[2])]$, i.e. $\langle \lstiC{POP x;}\invSC{(POP x)}, \sigma \rangle \Downarrow \sigma[\lstiC{x} \to (0,[2])]$, but $\sigma[x \to (0,[2])] \neq \sigma$, which means that the sequential composition of ``\lstiC{POP x}'', followed by its inverse ``\lstiC{PUSH x}'', starting from $\sigma$, does not lead back to $\sigma$.

The problem is that ``\lstiC{POP x}'' overwrites the current value of \lstiC{x}. Therefore, when trying to reverse the operation by applying ``\lstiC{PUSH x}'', the value to which \lstiC{x} must be reset is unknown. The \nsemantics will simply always set \lstiC{x} to the default value $0$. This is certainly not surprising, remarking that if we want be coherent with \TIF perspective, the operational semantics of \SCORE must be more sophisticated than \nsemantics.